\section{Ensemble and stochastic intermediate models}
\label{sec:ensembles}

Section~\ref{sec:upscaling} treats the intermediate model as one deterministic function to be smoothed. A second
line of work keeps \emph{many} models, or a random one, and lets the optimizer see the spread between them. Three
uses can be separated: averaging to reduce error, randomness to drive exploration, and stochastic gradients that let
first-order methods leave poor basins. They are rarely studied together, and almost never as a choice of
optimizer-facing surrogate.

\paragraph{Averaging ensembles.}
In engineering design the ensemble of surrogates is a weighted average of different model classes, with weights
from cross-validation error: the PRESS-weighted average surrogate \cite{goel2007ensemble}, weights optimised to
minimise the ensemble's cross-validation error \cite{acar2009ensemble}, and the finding that most of the benefit comes
from generating many diverse surrogates and selecting by PRESS, with weighting gains that shrink in high dimension
\cite{viana2009multiple}. A recent review \cite{chen2024ensemble} reports ensembles as the usual remedy for not knowing
which surrogate is best, which is the selection problem of this paper seen from the averaging side. In machine
learning the ensemble is homogeneous and randomised: bagging fits one class to bootstrap samples
\cite{breiman1996bagging}, and subsampling without replacement (subagging) has the same variance-reduction
mechanism with a cleaner analysis \cite{buhlmann2002analyzing}; random forests add feature subsampling
\cite{breiman2001random}; deep ensembles train networks from different initialisations
\cite{lakshminarayanan2017simple}; Monte-Carlo dropout samples sub-networks \cite{gal2016dropout}; and randomised
prior functions add a fixed random network to each member so that uncertainty exists away from the data, which
plain bootstrapping lacks \cite{osband2018randomized}. All of these give a mean predictor and a spread; the spread is
what acquisition rules (Section~\ref{sec:uq}) consume.

\paragraph{Randomness for exploration.}
The member-to-member variation can replace a posterior. Thompson sampling optimises one random draw of the model per
step; with ensembles the draw is one member, and the same idea underlies trust-region Thompson sampling in TuRBO
\cite{eriksson2019scalable}. Multiple-surrogate EGO adds one point per surrogate per cycle and imports uncertainty
estimates from one surrogate to another, so that models without native uncertainty can be used
\cite{viana2013efficient}; bootstrapped kriging replaces the plug-in kriging variance in expected improvement
\cite{kleijnen2012expected}; ensembles of Gaussian processes with different kernels are combined in a mixture
posterior for Thompson sampling \cite{lu2022surrogate}; and ensemble Bayesian optimization partitions the input space
with random Mondrian trees and fits a GP per cell \cite{wang2018batched}. Neural ensembles are a common alternative
to GPs in Bayesian optimization \cite{springenberg2016bayesian}, but a systematic comparison found deep ensembles
explore poorly and level off below other surrogates on standard benchmarks \cite{li2024study}. Tree ensembles are
the surrogate of SMAC \cite{hutter2011sequential}. In evolutionary computation, surrogate-assisted memetic algorithms
use ensembles and smoothed surrogates inside local search \cite{lim2010generalizing,jin2011surrogate}, and the
stochastic RBF method chooses each new point from random candidates scored by RBF prediction and distance, with
probabilistic convergence to the global minimum \cite{regis2007stochastic}. For noisy simulators, stochastic
kriging separates simulation noise from response-surface uncertainty \cite{ankenman2010stochastic}, which is the
regime where an averaged intermediate model is most defensible.

\paragraph{Stochastic gradients and smoothing.}
A deterministic gradient method stops in the first basin. Two mechanisms change that without a new surrogate. First,
noise in the gradient: Langevin dynamics adds Gaussian noise to stochastic gradient steps and has non-asymptotic
guarantees of reaching approximate global minimisers of non-convex objectives \cite{raginsky2017nonconvex}; using one
random ensemble member's gradient per step is a model-uncertainty version of the same noise. Second, smoothing:
convolving a non-smooth function with a Gaussian or uniform kernel gives a smooth surrogate whose gradient can be
estimated from random perturbations \cite{nesterov2017random,duchi2012randomized}. For a piecewise-constant tree
ensemble this is the only way to obtain useful gradient information, and a smoothing upscaler
(Section~\ref{sec:upscaling}) is its deterministic counterpart. Particle methods take the ensemble one step further
and make the optimizer itself an ensemble: consensus-based optimization moves particles towards a weighted consensus
point with an exploration noise \cite{pinnau2017consensus} and converges globally under mild conditions
\cite{fornasier2024consensus}, and ensemble Kalman inversion is a derivative-free ensemble optimizer from inverse
problems \cite{iglesias2013ensemble}.

\paragraph{Sums of local models.}
A surrogate that is a sum of simple local pieces is both smooth and naturally stochastic. The partition-of-unity
method blends local approximations with weights that sum to one \cite{babuska1997partition}; SNOBFIT fits local
quadratic models around the current best points and branches the domain to keep a global search
\cite{huyer2008snobfit}. Drawing one local piece by its weight gives a stochastic gradient whose expectation is the
weighted local gradient. We found no study that uses such a decomposition to run stochastic first-order methods on a
fitted surrogate.

\paragraph{Optimising tree ensembles directly.}
Instead of smoothing a tree ensemble, it can be optimised exactly: a tree ensemble is piecewise constant on a finite
partition, so its global minimum over a box is a mixed-integer linear program \cite{misic2020optimization}. ENTMOOT
solves this with a distance-based uncertainty term for Bayesian optimization over gradient-boosted trees
\cite{thebelt2021entmoot}, and tree-ensemble kernels give trees a GP-like posterior whose acquisition is a
mixed-integer second-order cone program \cite{thebelt2022tree}. This is the strongest alternative to upscaling for the
thesis' question of whether an accurate tree model can be used for design optimization.

\paragraph{State of the art and gaps.}
Ensembles are well established for accuracy (averaging) and for exploration (Thompson sampling, multiple-surrogate
EGO, SMAC), and tree ensembles can be optimised exactly. What is missing is a matched comparison of these routes from
the optimizer's side: (i) averaging versus sampling members versus smoothing versus exact mixed-integer optimization
for the same fitted ensemble, scored by regret on the true function rather than by accuracy; (ii) whether the
gradient noise from sampling members helps stochastic first-order methods find better optima than the deterministic
mean gradient, and at what cost in evaluations; (iii) whether the answer depends on the landscape type, as the
accuracy results of this paper suggest it will.

This study covers part of that space. Averaging ensembles (PRESS-weighted average, deep ensembles, Monte-Carlo
dropout, random and quantile forests, GBT committees), surrogate-managed loops with randomised acquisition (TuRBO,
DYCORS), and two stochastic upscalers with a stochastic-gradient arm (a partition of unity of local quadratics, and a
subsample ensemble of smooth members) are included. Methods the literature suggests adding, in order of how directly
they answer (i)--(iii): exact mixed-integer optimization of tree ensembles \cite{misic2020optimization,thebelt2021entmoot};
Thompson sampling over ensemble members as a fit-once optimizer (optimise one random member per restart);
Gaussian-smoothing gradient estimates on any surrogate \cite{nesterov2017random}; Langevin-type noise on the mean
gradient as the control for member sampling \cite{raginsky2017nonconvex}; consensus-based optimization as an ensemble
optimizer \cite{pinnau2017consensus}; optimised-weight ensembles \cite{acar2009ensemble} and randomised-prior members
\cite{osband2018randomized} as intermediate models; and bootstrapped-kriging expected improvement in the managed loops
\cite{kleijnen2012expected}.
